\section{Related Work}

Chapter 2 explained why the project uses a lightweight peer-to-peer discovery
layer. This chapter has a different purpose: it positions DIAL relative to
existing systems and papers. The main point is that DIAL combines two families
of ideas that are often studied separately: peer-to-peer discovery and LLM
routing.

\subsection{Peer-to-Peer Discovery and Decentralized Serving}

The networking part of DIAL follows the peer-to-peer tradition. Kademlia
provides the structured DHT idea used to locate providers by key
\citep{maymounkov2002kademlia}, while libp2p provides practical building
blocks for peer identity, transports, protocols, and DHT participation
\citep{libp2p}. BitTorrent and IPFS are useful background because they show
that open networks can coordinate resource exchange and content retrieval
without relying on a central content server
\citep{cohen2003bittorrent,benet2014ipfs}.

DIAL applies this pattern to LLM services rather than files or content blocks.
Each peer is treated as a service with a profile and a model backend. This
differs from systems such as Petals, which distribute parts of one large model
across many machines \citep{borzunov2023petals}. It is closer in spirit to
decentralized serving overlays such as PlanetServe, but DIAL studies a smaller
question in a local prototype: whether peers can advertise capabilities,
discover each other, and route queries measurably \citep{fang2026planetserve}.
Work such as Lattica is also relevant because open AI networks eventually need robust
connectivity across difficult network conditions, but that is outside the
evaluated scope here \citep{yang2025lattica}.

\subsection{LLM Routing, User Needs, and Model Profiles}

The routing part of DIAL connects to work on choosing between models. FrugalGPT
and RouteLLM show that model selection can trade off quality, cost, and
latency instead of sending every prompt to the same model
\citep{chen2023frugalgpt,ong2025routellm}. Surveys of routing and cascading
place these systems in a wider family of decision layers for LLM applications
\citep{moslem2026dynamicrouting}.

A related question is what a user request should be classified into before it
can be routed. TUNA studies real human-AI conversations and organizes user
needs into broad modes such as information seeking, synthesis, procedural
guidance, content creation, social interaction, and meta-conversation
\citep{shelby2025tuna}. Super-NaturalInstructions and HELM are useful from the
task and evaluation side: they show that language-model behavior is commonly
studied through families of tasks rather than one single capability
\citep{wang2022supernaturalinstructions,liang2022helm}. These works do not
define DIAL's exact eight labels, but they support the choice to use a small,
interpretable set of common request families for a prototype router.

For DIAL, the most important lesson is that routing needs profiles. Universal
Model Routing, RouteProfile, and vLLM Semantic Router all represent models or
requests with richer signals than a model name alone
\citep{jitkrittum2025uniroute,xu2026routeprofile,vllmsemanticrouter2026}. DIAL
uses a simpler capability vector because the goal is not to design the best
semantic router. The goal is to make profile-based routing work across peers
that are discovered dynamically through the network.

\subsection{Agent Discovery and PPAI}

PPAI is the closest conceptual neighbor because it describes personalized edge
LLM agents that can be matched to user demand through query-agent scoring and
load-aware delegation \citep{wang2026ppai}. DIAL does not implement the full
PPAI matching algorithm. Instead, it implements the lower layer that such a
system would need: peers join, publish capabilities, exchange profiles, route a
query, and expose traces that can be inspected.

Recent agent-discovery work makes the same infrastructure problem explicit.
Agents need to announce capabilities, be found, be ranked, and remain usable
as peers join, leave, or fail \citep{guo2026agentdiscovery}. Usable agent
discovery also notes that node availability and service readiness are related
but distinct problems
\citep{dazzi2026usableagentdiscovery}. For DIAL, this means discovery is only
the first step: the router still needs profile and availability information to
choose between candidates.

\subsection{Future Layers: Aggregation, Privacy, and Trust}

The current system selects one node to answer a query. The papers below are
useful mainly because they show what could be layered on top of that first
routing step. LLM-Blender and Mixture-of-Agents show that multiple model
outputs can be ranked, fused, or layered to improve quality in multi-model settings
\citep{jiang2023llmblender,wang2024mixtureofagents}. LLM-as-a-Verifier is
relevant if future versions need a separate model to judge routed answers
\citep{kwok2026llmverifier}. Privacy-preserving routing, verifiable inference,
and blockchain-based AI trust work point to mechanisms that become relevant if
the network moves beyond controlled local experiments and into open deployments
\citep{wu2026privacyrouting,wang2026verillm,nassar2020blockchain}.

These papers do not describe what DIAL already implements. They clarify the
layers that would be needed after discovery and routing: answer aggregation,
privacy policies, stronger identity, reputation, and verification.
